%% Layer 1 in detail: three stages, the artefact that flows between them, and where the
%% two measurement instruments attach.
%%
%% Hand-authored TikZ (not generated): it depicts a design, so there is no data source to
%% derive it from. The instrument annotations carry as much weight as the stages -- they are
%% what makes reading and structuring separately accountable, and therefore what makes the
%% attribution result in this thesis possible at all.
%%
%% Artefacts ride on the arrows rather than occupying their own boxes: the stages are what
%% the reader must remember, and five boxes in a row read faster than nine.
%%
%% Input from a chapter with:  %% Layer 1 in detail: three stages, the artefact that flows between them, and where the
%% two measurement instruments attach.
%%
%% Hand-authored TikZ (not generated): it depicts a design, so there is no data source to
%% derive it from. The instrument annotations carry as much weight as the stages -- they are
%% what makes reading and structuring separately accountable, and therefore what makes the
%% attribution result in this thesis possible at all.
%%
%% Artefacts ride on the arrows rather than occupying their own boxes: the stages are what
%% the reader must remember, and five boxes in a row read faster than nine.
%%
%% Input from a chapter with:  %% Layer 1 in detail: three stages, the artefact that flows between them, and where the
%% two measurement instruments attach.
%%
%% Hand-authored TikZ (not generated): it depicts a design, so there is no data source to
%% derive it from. The instrument annotations carry as much weight as the stages -- they are
%% what makes reading and structuring separately accountable, and therefore what makes the
%% attribution result in this thesis possible at all.
%%
%% Artefacts ride on the arrows rather than occupying their own boxes: the stages are what
%% the reader must remember, and five boxes in a row read faster than nine.
%%
%% Input from a chapter with:  %% Layer 1 in detail: three stages, the artefact that flows between them, and where the
%% two measurement instruments attach.
%%
%% Hand-authored TikZ (not generated): it depicts a design, so there is no data source to
%% derive it from. The instrument annotations carry as much weight as the stages -- they are
%% what makes reading and structuring separately accountable, and therefore what makes the
%% attribution result in this thesis possible at all.
%%
%% Artefacts ride on the arrows rather than occupying their own boxes: the stages are what
%% the reader must remember, and five boxes in a row read faster than nine.
%%
%% Input from a chapter with:  \input{figures/pipeline.tex}

\begin{tikzpicture}[
    font=\small,
    node distance=9mm,
    stage/.style={
      draw=horusaccent, line width=0.8pt, rounded corners=2pt,
      fill=horuswash, text width=2.6cm, align=center,
      minimum height=1.6cm, inner sep=4pt,
    },
    endpoint/.style={
      draw=horusmuted, line width=0.6pt, rounded corners=1pt,
      fill=white, text width=2.0cm, align=center,
      font=\scriptsize, minimum height=1.0cm, inner sep=4pt,
    },
    instrument/.style={
      draw=horuswarn, line width=0.7pt, dashed, rounded corners=2pt,
      fill=white, text=horuswarn, text width=3.4cm, align=center,
      font=\scriptsize, inner sep=4pt,
    },
    flow/.style={-{Stealth[length=2.4mm]}, line width=0.8pt, draw=horusaccent},
    probe/.style={-{Stealth[length=2.0mm]}, line width=0.6pt, draw=horuswarn, dashed},
    onflow/.style={font=\scriptsize, text=horusink, align=center, inner sep=2pt},
  ]

  \node[endpoint] (page) {page image\\\itshape 300\,dpi};
  \node[stage, right=of page] (reader) {\textbf{Reader}\\[2pt]\scriptsize a vision model
    transcribes the page};
  \node[stage, right=22mm of reader] (structurer) {\textbf{Structurer}\\[2pt]\scriptsize a
    language model fills the schema};
  \node[stage, right=18mm of structurer] (repair) {\textbf{Validate\\and repair}\\[2pt]
    \scriptsize deterministic rules, no model};
  \node[endpoint, right=of repair] (final) {validated\\record};

  \draw[flow] (page) -- (reader);
  \draw[flow] (reader) -- node[onflow, above=1pt] {transcript\\\itshape plain text}
    (structurer);
  \draw[flow] (structurer) -- node[onflow, above=1pt] {draft record\\\itshape JSON} (repair);
  \draw[flow] (repair) -- (final);

  \begin{scope}[on background layer]
    \node[
      draw=horusmuted, line width=0.5pt, rounded corners=3pt, dotted,
      fit=(page)(reader)(structurer)(repair)(final),
      inner xsep=7pt, inner ysep=9pt,
    ] (box) {};
  \end{scope}

  %% Both measurement instruments act on the transcript, which is why they share one box and
  %% one attachment point. Instrument 1 removes the reading stage from the measurement without
  %% changing anything else, making the structurer's own ceiling observable. Instrument 2 asks,
  %% for each individual miss, whether the value was in the transcript at all -- the question
  %% that assigns an error to a stage.
  \coordinate (tap) at ($(reader.east)!0.5!(structurer.west)$);
  \node[
    draw=horuswarn, line width=0.7pt, dashed, rounded corners=2pt,
    fill=white, text=horuswarn, text width=7.4cm, align=left,
    font=\scriptsize, inner sep=5pt,
    below=13mm of box.south, anchor=north,
  ] (instruments) {%
    \textbf{The transcript is where both measurement instruments attach:}\\[2pt]
    \textbf{1.} replace it with text rendered from the ground truth --- this removes the
    reading stage from the measurement and changes nothing else\\[1pt]
    \textbf{2.} for each individual miss, ask whether the value was present in it at all%
  };
  \draw[probe] (instruments.north) -- (tap);

\end{tikzpicture}


\begin{tikzpicture}[
    font=\small,
    node distance=9mm,
    stage/.style={
      draw=horusaccent, line width=0.8pt, rounded corners=2pt,
      fill=horuswash, text width=2.6cm, align=center,
      minimum height=1.6cm, inner sep=4pt,
    },
    endpoint/.style={
      draw=horusmuted, line width=0.6pt, rounded corners=1pt,
      fill=white, text width=2.0cm, align=center,
      font=\scriptsize, minimum height=1.0cm, inner sep=4pt,
    },
    instrument/.style={
      draw=horuswarn, line width=0.7pt, dashed, rounded corners=2pt,
      fill=white, text=horuswarn, text width=3.4cm, align=center,
      font=\scriptsize, inner sep=4pt,
    },
    flow/.style={-{Stealth[length=2.4mm]}, line width=0.8pt, draw=horusaccent},
    probe/.style={-{Stealth[length=2.0mm]}, line width=0.6pt, draw=horuswarn, dashed},
    onflow/.style={font=\scriptsize, text=horusink, align=center, inner sep=2pt},
  ]

  \node[endpoint] (page) {page image\\\itshape 300\,dpi};
  \node[stage, right=of page] (reader) {\textbf{Reader}\\[2pt]\scriptsize a vision model
    transcribes the page};
  \node[stage, right=22mm of reader] (structurer) {\textbf{Structurer}\\[2pt]\scriptsize a
    language model fills the schema};
  \node[stage, right=18mm of structurer] (repair) {\textbf{Validate\\and repair}\\[2pt]
    \scriptsize deterministic rules, no model};
  \node[endpoint, right=of repair] (final) {validated\\record};

  \draw[flow] (page) -- (reader);
  \draw[flow] (reader) -- node[onflow, above=1pt] {transcript\\\itshape plain text}
    (structurer);
  \draw[flow] (structurer) -- node[onflow, above=1pt] {draft record\\\itshape JSON} (repair);
  \draw[flow] (repair) -- (final);

  \begin{scope}[on background layer]
    \node[
      draw=horusmuted, line width=0.5pt, rounded corners=3pt, dotted,
      fit=(page)(reader)(structurer)(repair)(final),
      inner xsep=7pt, inner ysep=9pt,
    ] (box) {};
  \end{scope}

  %% Both measurement instruments act on the transcript, which is why they share one box and
  %% one attachment point. Instrument 1 removes the reading stage from the measurement without
  %% changing anything else, making the structurer's own ceiling observable. Instrument 2 asks,
  %% for each individual miss, whether the value was in the transcript at all -- the question
  %% that assigns an error to a stage.
  \coordinate (tap) at ($(reader.east)!0.5!(structurer.west)$);
  \node[
    draw=horuswarn, line width=0.7pt, dashed, rounded corners=2pt,
    fill=white, text=horuswarn, text width=7.4cm, align=left,
    font=\scriptsize, inner sep=5pt,
    below=13mm of box.south, anchor=north,
  ] (instruments) {%
    \textbf{The transcript is where both measurement instruments attach:}\\[2pt]
    \textbf{1.} replace it with text rendered from the ground truth --- this removes the
    reading stage from the measurement and changes nothing else\\[1pt]
    \textbf{2.} for each individual miss, ask whether the value was present in it at all%
  };
  \draw[probe] (instruments.north) -- (tap);

\end{tikzpicture}


\begin{tikzpicture}[
    font=\small,
    node distance=9mm,
    stage/.style={
      draw=horusaccent, line width=0.8pt, rounded corners=2pt,
      fill=horuswash, text width=2.6cm, align=center,
      minimum height=1.6cm, inner sep=4pt,
    },
    endpoint/.style={
      draw=horusmuted, line width=0.6pt, rounded corners=1pt,
      fill=white, text width=2.0cm, align=center,
      font=\scriptsize, minimum height=1.0cm, inner sep=4pt,
    },
    instrument/.style={
      draw=horuswarn, line width=0.7pt, dashed, rounded corners=2pt,
      fill=white, text=horuswarn, text width=3.4cm, align=center,
      font=\scriptsize, inner sep=4pt,
    },
    flow/.style={-{Stealth[length=2.4mm]}, line width=0.8pt, draw=horusaccent},
    probe/.style={-{Stealth[length=2.0mm]}, line width=0.6pt, draw=horuswarn, dashed},
    onflow/.style={font=\scriptsize, text=horusink, align=center, inner sep=2pt},
  ]

  \node[endpoint] (page) {page image\\\itshape 300\,dpi};
  \node[stage, right=of page] (reader) {\textbf{Reader}\\[2pt]\scriptsize a vision model
    transcribes the page};
  \node[stage, right=22mm of reader] (structurer) {\textbf{Structurer}\\[2pt]\scriptsize a
    language model fills the schema};
  \node[stage, right=18mm of structurer] (repair) {\textbf{Validate\\and repair}\\[2pt]
    \scriptsize deterministic rules, no model};
  \node[endpoint, right=of repair] (final) {validated\\record};

  \draw[flow] (page) -- (reader);
  \draw[flow] (reader) -- node[onflow, above=1pt] {transcript\\\itshape plain text}
    (structurer);
  \draw[flow] (structurer) -- node[onflow, above=1pt] {draft record\\\itshape JSON} (repair);
  \draw[flow] (repair) -- (final);

  \begin{scope}[on background layer]
    \node[
      draw=horusmuted, line width=0.5pt, rounded corners=3pt, dotted,
      fit=(page)(reader)(structurer)(repair)(final),
      inner xsep=7pt, inner ysep=9pt,
    ] (box) {};
  \end{scope}

  %% Both measurement instruments act on the transcript, which is why they share one box and
  %% one attachment point. Instrument 1 removes the reading stage from the measurement without
  %% changing anything else, making the structurer's own ceiling observable. Instrument 2 asks,
  %% for each individual miss, whether the value was in the transcript at all -- the question
  %% that assigns an error to a stage.
  \coordinate (tap) at ($(reader.east)!0.5!(structurer.west)$);
  \node[
    draw=horuswarn, line width=0.7pt, dashed, rounded corners=2pt,
    fill=white, text=horuswarn, text width=7.4cm, align=left,
    font=\scriptsize, inner sep=5pt,
    below=13mm of box.south, anchor=north,
  ] (instruments) {%
    \textbf{The transcript is where both measurement instruments attach:}\\[2pt]
    \textbf{1.} replace it with text rendered from the ground truth --- this removes the
    reading stage from the measurement and changes nothing else\\[1pt]
    \textbf{2.} for each individual miss, ask whether the value was present in it at all%
  };
  \draw[probe] (instruments.north) -- (tap);

\end{tikzpicture}


\begin{tikzpicture}[
    font=\small,
    node distance=9mm,
    stage/.style={
      draw=horusaccent, line width=0.8pt, rounded corners=2pt,
      fill=horuswash, text width=2.6cm, align=center,
      minimum height=1.6cm, inner sep=4pt,
    },
    endpoint/.style={
      draw=horusmuted, line width=0.6pt, rounded corners=1pt,
      fill=white, text width=2.0cm, align=center,
      font=\scriptsize, minimum height=1.0cm, inner sep=4pt,
    },
    instrument/.style={
      draw=horuswarn, line width=0.7pt, dashed, rounded corners=2pt,
      fill=white, text=horuswarn, text width=3.4cm, align=center,
      font=\scriptsize, inner sep=4pt,
    },
    flow/.style={-{Stealth[length=2.4mm]}, line width=0.8pt, draw=horusaccent},
    probe/.style={-{Stealth[length=2.0mm]}, line width=0.6pt, draw=horuswarn, dashed},
    onflow/.style={font=\scriptsize, text=horusink, align=center, inner sep=2pt},
  ]

  \node[endpoint] (page) {page image\\\itshape 300\,dpi};
  \node[stage, right=of page] (reader) {\textbf{Reader}\\[2pt]\scriptsize a vision model
    transcribes the page};
  \node[stage, right=22mm of reader] (structurer) {\textbf{Structurer}\\[2pt]\scriptsize a
    language model fills the schema};
  \node[stage, right=18mm of structurer] (repair) {\textbf{Validate\\and repair}\\[2pt]
    \scriptsize deterministic rules, no model};
  \node[endpoint, right=of repair] (final) {validated\\record};

  \draw[flow] (page) -- (reader);
  \draw[flow] (reader) -- node[onflow, above=1pt] {transcript\\\itshape plain text}
    (structurer);
  \draw[flow] (structurer) -- node[onflow, above=1pt] {draft record\\\itshape JSON} (repair);
  \draw[flow] (repair) -- (final);

  \begin{scope}[on background layer]
    \node[
      draw=horusmuted, line width=0.5pt, rounded corners=3pt, dotted,
      fit=(page)(reader)(structurer)(repair)(final),
      inner xsep=7pt, inner ysep=9pt,
    ] (box) {};
  \end{scope}

  %% Both measurement instruments act on the transcript, which is why they share one box and
  %% one attachment point. Instrument 1 removes the reading stage from the measurement without
  %% changing anything else, making the structurer's own ceiling observable. Instrument 2 asks,
  %% for each individual miss, whether the value was in the transcript at all -- the question
  %% that assigns an error to a stage.
  \coordinate (tap) at ($(reader.east)!0.5!(structurer.west)$);
  \node[
    draw=horuswarn, line width=0.7pt, dashed, rounded corners=2pt,
    fill=white, text=horuswarn, text width=7.4cm, align=left,
    font=\scriptsize, inner sep=5pt,
    below=13mm of box.south, anchor=north,
  ] (instruments) {%
    \textbf{The transcript is where both measurement instruments attach:}\\[2pt]
    \textbf{1.} replace it with text rendered from the ground truth --- this removes the
    reading stage from the measurement and changes nothing else\\[1pt]
    \textbf{2.} for each individual miss, ask whether the value was present in it at all%
  };
  \draw[probe] (instruments.north) -- (tap);

\end{tikzpicture}
